\documentclass[11pt]{article}
\usepackage[margin=1in]{geometry}
\usepackage[T1]{fontenc}
\usepackage[utf8]{inputenc}
\usepackage{lmodern}
\usepackage{graphicx}
\usepackage{booktabs}
\usepackage{hyperref}
\usepackage{float}
\title{Non-LLM Complexity Addendum}
\author{}
\date{April 7, 2026}
\begin{document}
\maketitle
\section*{Scope}
This addendum extends the approved-solver complexity analysis to the non-LLM systems stored in this repository.
\section*{Data}
\begin{itemize}
  \item ARC-1 approved eval overlap: 38 tasks
  \item ARC-2 approved eval overlap: 18 tasks
  \item ARC-2 tasks with both human and non-LLM difficulty: 17 tasks
  \item ARC-1 non-LLM profiles: 11
  \item ARC-2 non-LLM profiles: 28
\end{itemize}
\section*{Methods}
For each dataset I built a task-by-profile success matrix and derived smoothed-logit, PCA, and penalized 1PL item-difficulty summaries. The key-test table uses the smoothed-logit difficulty as the primary outcome. Statistical inference used bootstrap confidence intervals, permutation p-values for direct correlations, bootstrap difference-of-correlation tests, and BH-FDR correction across the 9 key hypotheses.
\section*{Headline Results}
Best ARC-2 structural metric: Gzip bytes ($r=0.504$, $q=0.217$).\\
Best ARC-1 structural metric: Gzip bytes ($r=0.366$, $q=0.191$).\\
Best pooled metric after within-dataset standardization: Gzip bytes ($r=0.378$, $q=0.062$).
\begin{center}
\small
\begin{tabular}{lrrrrrr}
\toprule
ID & Estimate & CI low & CI high & p & q & FDR \\
\midrule
N1 & 0.113 & -0.364 & 0.518 & 0.657 & 0.739 & no \\
N2 & 0.062 & -0.381 & 0.440 & 0.818 & 0.818 & no \\
N3 & 0.469 & 0.109 & 0.864 & 0.518 & 0.721 & no \\
N4 & 0.352 & -0.157 & 0.660 & 0.148 & 0.667 & no \\
N5 & 0.167 & -0.170 & 0.742 & 0.561 & 0.721 & no \\
N6 & 0.346 & -0.044 & 0.658 & 0.508 & 0.721 & no \\
N7 & 0.356 & -0.066 & 0.817 & 0.495 & 0.721 & no \\
N8 & 0.255 & -0.042 & 0.507 & 0.125 & 0.667 & no \\
N9 & 0.113 & -0.211 & 0.399 & 0.499 & 0.721 & no \\
\bottomrule
\end{tabular}
\end{center}
\section*{Interpretation}
The cleanest read from this first pass is that non-LLM item difficulty is not random with respect to approved solver complexity, but the structural-complexity link is weaker and less focused than the LLM-side result. The ARC-2 non-LLM profiles still line up more naturally with the LLM difficulty axis than with the human axis, and once LLM difficulty is controlled the remaining structural signal in the non-LLM axis is smaller and less secure. A pooled within-dataset-standardized view strengthens the directional pattern, which suggests that low power is part of the story in the split analyses.
\begin{figure}[H]
  \centering
  \includegraphics[width=\linewidth]{chart_non_llm_arc2_relationships.png}
  \caption{ARC-2 non-LLM relationships.}
\end{figure}
\begin{figure}[H]
  \centering
  \includegraphics[width=\linewidth]{chart_non_llm_complexity_comparison.png}
  \caption{Human vs LLM vs non-LLM complexity comparisons on the shared 17-task ARC-2 overlap.}
\end{figure}
\begin{figure}[H]
  \centering
  \includegraphics[width=\linewidth]{chart_non_llm_arc1_sidecar.png}
  \caption{ARC-1 sidecar results.}
\end{figure}
\begin{figure}[H]
  \centering
  \includegraphics[width=0.8\linewidth]{chart_non_llm_trm_trajectory.png}
  \caption{Exploratory TRM trajectory vs complexity.}
\end{figure}
\end{document}
